\subsection{The tensor defect in an actual repaired partial source}

\begin{proposition}[A-L17: Proved; coefficient base change of AT23]
\label{AL17}
Let $Y$ be the integral strict simplicial PD algebra of
AT23 for a prime $p$ and a weight $w\geq3$, and let
$S=\Gamma_{\mathbb Z}(T)$. In every simplex $n$ its
coefficient map induces $H^0(Y_n)=S$ and
$H^i(Y_n)=0$ for $1\leq i\leq3$.
The actual strict coefficient base change
\[
 Y\longrightarrow Y\otimes_S(S/2S)=Y\otimes_{\mathbb Z}\mathbb F_2
 \tag{AL.91}
\]
has, in simplex two,
\[
 H^3(Y_2\otimes_S S/2S)
      \cong (S/2S)\langle2w\rangle,
 \qquad [v x']\longleftrightarrow1,
 \tag{AL.92}
\]
where $x'=x-pu$ is the actual integral generator of
AT23. The class is nonzero in every weight at least
$2w$ after multiplying by the corresponding coefficient.
These are partial sources with the properties proved
in AT23; no assertion of eventual vertical admissibility
is added here.
\end{proposition}

\begin{proof}
The integral inverse change (AT.147) identifies $Y_n$
with $G_n$ tensored with one even identity disk and
one odd identity disk for each basis element of
$D(3)_n$. This is a finite set in each simplex.
The actual $S$-linear coefficient contraction of $G_n$
and the integral monomial contractions of all even
disks identify its vertical cohomology, before and
after the base change (AL.91), with that of the
tensor product of the odd disks
\[
 \Lambda(v_j)\otimes\Gamma(x'_j),\qquad
                 dv_j=-x'_j,\quad dx'_j=0.
 \tag{AL.93}
\]
All maps in these contractions are integral and
$S$-linear, so their tensor products remain the same
contractions after the specified coefficient quotient.

Decompose (AL.93) by its vector of arities, one arity
for each $j$. Its factor of positive arity $a$ is
the actual two-term complex
\[
 S\,v_j\gamma_{a-1}(x'_j)
       \xrightarrow{\ -a\ }S\,\gamma_a(x'_j),
 \quad\text{in degrees }2a-1,2a.
 \tag{AL.94}
\]
A factor of arity one is contractible integrally,
and so is any tensor summand containing such a
factor. A remaining nonconstant factor has arity
at least two and starts in form degree at least
three. Two such factors start in degree at least
six. Thus the only nonconstant summands meeting
degrees at most three are single factors of arity
two, and their degree-three differential is the
injective map $-2:S\to S$. This proves coefficient
cohomology through degree three in every simplex.

In simplex two, $D(3)_2$ has one basis element.
After (AL.91), the arity-two differential in (AL.94)
is zero; it has no incoming group in degree two.
Arity-one summands remain contractible and all higher
arities begin above degree three. This proves the
entire group (AL.92), with the displayed generator.
Each $S/2S$ weight component is the nonzero group
$\mathbb F_2 T^{[a]}$, so all the asserted coefficient
multiples remain nonzero.

The coefficient inclusion and any vertex evaluation
give the actual splitting required by A-L15 for
$M=Y_n$. All its cochain terms are weight-graded
$S$-free with nonnegative generator weights. Taking
$d_0=3$ and $N=S/2S$ in A-L16 identifies (AL.92)
with the precise group
$\operatorname{Tor}^S_1(Y_2^4/B^4(Y_2),S/2S)$.
This is the stated realization of that tensor defect
by the actual partial-source construction and map.
\end{proof}
